\chapter{Provider-Free Reproducibility Guide}

\section{Purpose}

This appendix describes how a reviewer can validate the final evidence without spending provider credit. It is not an instruction to rerun paid experiments. The review begins from the canonical source-corrected data and frozen release/evidence packages already contained in the repository.

\section{Recommended sequence}

\begin{enumerate}[leftmargin=1.6em]
\item Confirm the repository revision and inspect the working tree for unexpected changes.
\item Verify the canonical source-corrected data hash and the 1,568-record/zero-unresolved accounting.
\item Verify the source-corrected full-population analysis artifact identity and the result ledger.
\item Run the corrected-evidence and release-v4 verification tools.
\item Restore the release-v4 database into a disposable PostgreSQL instance.
\item Start the backend against that disposable database only and call the controlled-results endpoint.
\item Reconcile returned RQ1--RQ6 values and the RQ7 primary/secondary sections with the stored authoritative analysis runs.
\item Run backend tests, frontend tests, linting, TypeScript checks, and the production build.
\end{enumerate}

\section{What a successful check establishes}

Successful verification establishes that the checked artifact and database snapshot have the expected bytes, that the backend can select the intended analysis records, and that the user interface can consume the API. It does not establish that a provider would behave identically today, because the workflow purposefully makes no provider request. This distinction is central: scientific reproducibility of stored evidence is different from repeat execution of a live model service.

\section{Failure interpretation}

If a checksum fails, the appropriate response is to stop and determine whether the artifact changed. If an analysis run is missing or duplicate, the controlled-results API should fail closed. If a historical explorer has missing answer provenance, it should return unavailable provenance rather than fabricate text. If a live sandbox request is disabled, it should be rejected before provider transport. These failures are protective signals, not reasons to replace stored values with constants.

\section{Preservation recommendations}

For archival use, preserve the repository revision, canonical source-corrected data, all evidence packages, release dumps, manifest hashes, and this report PDF. Preserve the distinction between historical Phase 11 evidence and newer source-corrected final evidence. When new research is performed, create a new named lineage; do not overwrite the identities summarized here.
